\documentclass[10pt,a4paper]{amsart}
\usepackage[margin=19mm]{geometry}
\usepackage{amsmath,amssymb,mathtools}
\usepackage[scr=rsfs,cal=euler]{mathalfa}
\usepackage{xfrac}
\usepackage[unicode,hidelinks,bookmarks=false]{hyperref}
\newcommand{\ie}{i.e.\ }
\newcommand{\cf}{cf.\ }
\newcommand{\resp}{respectively\ }
\newcommand{\Subsection}{\subsection}
\providecommand{\nobreakdash}{\nobreak\hbox{-}}
\setlength{\emergencystretch}{2em}
\multlinegap0pt
\usepackage{amssymb}
\newtheorem{lemma}{Lemma}
\newtheorem{proposition}{Proposition}
\newcommand{\hcf}{\operatorname{hcf}}
\title{A parallel wakeup problem and multi-room light switch strategies}
\author{John Haslegrave, Paul A. Russell, Mark Walters}
\date{Source: arXiv:2605.19488v1 (CC BY 4.0)}
\begin{document}
\maketitle

\noindent\emph{Source and licence.} A parallel wakeup problem and multi-room light switch strategies. Version arXiv:2605.19488v1 (CC BY 4.0), distributed by arXiv under the Creative Commons Attribution 4.0 International licence, http://creativecommons.org/licenses/by/4.0/, as recorded in the arXiv licence metadata for this exact version and shown on the abstract page https://arxiv.org/abs/2605.19488v1. Full licence notice: \texttt{licenses/arxiv-2605.19488-CC-BY-4.0.txt}.\par\smallskip

\noindent\textit{Editorial scope.} Counted unit: the article's proposition n-smaller (source0.tex lines 267--269). The article's complete original proof of it is delivered with it (source0.tex lines 270--289, one \texttt{proof} environment of 20 lines). No mathematics is written, repaired or re-derived: the statement and the proof are the article's own lines, in the article's own notation, and no retained line is substituted or re-flowed.  Retained uncounted as the source-local context the proof needs: lines 262--262.  Nothing was omitted from any retained span, so the package carries no omission record.  No retained line contains a citation, so the wrapper carries no bibliography and no bibliography entry.  No retained line contains a figure environment, an image-inclusion command or drawing code, so no image is delivered and the package declares no asset dependency.  Editorial formatting only, in addition: an AMS \texttt{amsart} preamble that restates the article's own declarations and macros used by the retained lines, namely the environments \texttt{lemma}, \texttt{proposition} on separate counters, together with the standard packages those lines need.  The retained environments print in this document as Lemma 1, Proposition 1 in order of appearance, whereas the article prints the same items with the numbering of its own full document.  The complete native source of the article as distributed in the arXiv e-print for this exact version is bundled unchanged as \texttt{sources/200918/source0.tex}.
	\begin{lemma}\label{leader-wins}Suppose that, for $n$ prisoners, $r$ rooms and $q$ states, the prisoners have a symmetric strategy ensuring that exactly one prisoner will at some point consider themselves a leader. Then $n$ prisoners have a symmetric winning strategy for $r$ rooms with $q+4$ states.\end{lemma}

	\begin{proposition}\label{n-smaller}
		For every $n$ there exists $q=q(n)$ such that for every $r>n$ with $\hcf(n,r)=1$, $n$ prisoners have a symmetric winning strategy for $r$ rooms with $q$ states.
	\end{proposition}

	\begin{proof}
		By Lemma \ref{leader-wins}, it is sufficient to show that the prisoners can find a leader using $q-4$ states. For simplicity of argument, we describe the leader-finding strategy using two switches in each room: a binary switch $S$ and an $n$-state switch $T$. These can be mapped onto states $0,\ldots, 2n-1$, giving $q=2n+4$.
		No prisoner will ever decrease the state of any $S$- or $T$-switch. When a prisoner first sees a room with the $S$-switch set to $0$, they move it to a $1$. They do this regardless of everything else. After this they never touch an $S$-switch again 
		Since $n < r$ this means that eventually $n$ rooms will be marked.
		
		We now use the following ``race'' strategy to find a leader. Each prisoner starts at level $n$. 
		On level $i$, each prisoner tries to move the $T$-switch in a certain number of rooms from any state that is at most $n-i$ to state $n-i+1$. 
		If $r/i$ is not an integer, prisoners on level $i$ attempt to do this in $\lceil r/i\rceil$ rooms. If $r/i$ is an integer, they attempt to do this in $\lceil n/i\rceil$
		rooms that have their $S$-switch set to $1$. Any prisoner who succeeds moves to level $i-1$.
		
		We argue by induction that at least one and at most $i$ prisoners reach level $i$. This is trivial
		for level $n$. Suppose this holds for $i > 1$. Since $\hcf(n, r) = 1$, one of $r/i$ and $n/i$ is not an integer,
		so either the target is $\lceil r/i\rceil$ and $i\lceil r/i\rceil>r$ or the target is $\lceil n/i\rceil$ marked rooms and $i\lceil n/i\rceil>n$; 
		in either case it is impossible for $i$ prisoners to reach the target. Suppose that no prisoners reach level $i-1$. By the induction
		hypothesis, at most $i$ prisoners (and at least one) reached level $i$. If the target is $\lceil r/i\rceil$, at least one of these prisoners must reach it,
		since otherwise the $T$-switches in at most $i(\lceil r/i\rceil-1) < r$ rooms reach state $n-i+1$ or above, and eventually one of the prisoners will be able to change another one.
		Similarly, if the target is $\lceil n/i\rceil$, eventually there will be $n$ marked rooms and thereafter at least one prisoner must reach the target from among these rooms.
		
		The unique prisoner who reaches level $1$ becomes the leader. 
	\end{proof}

\end{document}
